% BHU PYQ -- Manually transcribed & error-corrected
% Paper: BCHI-326 : Industrial Relations and Labour Laws
% Exam: B.Com. (Hons.) Semester VI Examination, 2022-23

\documentclass[11pt,a4paper]{article}
\usepackage[a4paper,top=2.5cm,bottom=2.5cm,left=2.8cm,right=2.8cm,headheight=14pt]{geometry}
\usepackage{amsmath,amssymb}
\usepackage[T1]{fontenc}
\usepackage[utf8]{inputenc}
\usepackage{lmodern,microtype,fancyhdr,enumitem,hyperref,lastpage,parskip}

\hypersetup{
    colorlinks=true,
    urlcolor=black,
    linkcolor=black,
    pdftitle={BCHI-326: Industrial Relations and Labour Laws -- BHU B.Com. (Hons.) Sem VI 2022-23}
}

\pagestyle{fancy}
\fancyhf{}
\renewcommand{\headrulewidth}{0.4pt}
\renewcommand{\footrulewidth}{0.4pt}
\fancyhead[L]{\small   \textit{Banaras Hindu University}}
\fancyhead[R]{\small   \textit{BCHI-326: Industrial Relations and Labour Laws}}
\fancyfoot[C]{\small Page  \thepage\ of \pageref{LastPage}}


\newlist{parts}{enumerate}{2}
\setlist[parts,1]{label=(\alph*),leftmargin=*,itemsep=5pt,topsep=3pt}
\setlist[parts,2]{label=(\roman*),leftmargin=*,itemsep=3pt,topsep=2pt}

\newcommand{\pts}[1]{\hfill   \textit{[#1]}}


\begin{document}


\begin{center}
  {\large\bfseries BANARAS HINDU UNIVERSITY}\\[2pt]
  {\normalsize B.Com. (Hons.) --- Semester VI Examination, 2022-23}\\[6pt]
  \rule{0.8\linewidth}{0.5pt}\\[6pt]
  {\large\bfseries Paper No. BCHI-326: Industrial Relations and Labour Laws}\\[4pt]
  {\normalsize Time: 3 Hours\hfill Maximum Marks: 70}
\end{center}

\rule{\linewidth}{0.4pt}

\smallskip



\noindent   \textit{Instructions: Attempt all questions. All questions carry equal marks (14 marks each).}

\bigskip




\noindent   \textbf{Question 1.}
What is meant by ``Industrial Relations''? Discuss the emerging trends in industrial relations in India.\pts{14}

\centerline{   \textbf{OR}}

Describe Dunlop's model of industrial relations along with the limitations of its applicability in the Indian context.\pts{14}

\medskip\rule{\linewidth}{0.2pt}




\noindent   \textbf{Question 2.}
Discuss the methods adopted by trade unions to achieve their objectives. Which of these methods will you suggest for Indian trade unions today?\pts{14}

\centerline{   \textbf{OR}}

Describe the various kinds of primary trade unions in India. Which of these is more prevalent in the country?\pts{14}

\medskip\rule{\linewidth}{0.2pt}




\noindent   \textbf{Question 3.}
Describe the statutory machinery for the settlement of industrial disputes in India. Which of these will you prefer for the country and why?\pts{14}

\centerline{   \textbf{OR}}

Explain the various causes of industrial disputes resulting in work stoppages with particular reference to India along with their relative importance.\pts{14}

\medskip\rule{\linewidth}{0.2pt}




\noindent   \textbf{Question 4.}
Give a brief description of the growth of collective bargaining in India along with the hurdles in its way.\pts{14}

\centerline{   \textbf{OR}}

What is meant by ``collective bargaining''? Briefly describe the important theories of collective bargaining.\pts{14}

\medskip\rule{\linewidth}{0.2pt}




\noindent   \textbf{Question 5.}
Give an assessment of the working of the Factories Act, 1948 and suggest measures for improving its effectiveness.\pts{14}

\centerline{   \textbf{OR}}

Describe the procedure laid down under the Minimum Wages Act, 1948 for the fixation and revision of current rates of wages. Suggest measures to make it more effective.\pts{14}

\vfill
\rule{\linewidth}{0.4pt}

\begin{center}\small
\href{https://bhu-exam.vercel.app/}{Click here to visit our website}\\[4pt]
\href{https://me-aryan.vercel.app/}{Click here to visit our contributor}\\[4pt]
\href{https://quashan-me.vercel.app/}{Click here to visit our contributor}
\end{center}

\end{document}
